% Supplementary Note 2: Study 1's robustness checks and the belief check. Numbers through macros only.
\suppnote{Study 1: robustness checks and the belief check}

This note gives the checks on Study~1's label penalty and on what Human-label participants wrote about machines, which the main text summarizes in its Results on Study~1.

Study~1's \SOneN{} participants rated the same ten AI-made images for creativity and for effort on 1-to-5 scales. A label told each participant that the artist, Alex, had made them either with design and sketching tools (the Human label) or by prompting an AI system (the AI label). The label's penalty is the mean rating under the Human label minus the mean under the AI label, so a positive penalty means that the AI label lowered the rating. Its primary estimate comes from a linear mixed model, a regression that gives each participant and each image its own baseline: \SOneHOneCreatDiff{} points on creativity and \SOneHOneEffDiff{} on effort, the registered confirmatory test of Study~1's first hypothesis (H1). All but one of the analyses below are named in Study~1's registration, as checks on that estimate or as counts to report, which makes them registered secondary analyses: reported, but not decisions about a hypothesis. The exception, the comparison with random groups, is exploratory description, as the registration names it. None replaces the primary estimate. Study~1 was registered on 8 October 2026, after its analyses had been run.

\paragraph*{Robustness of the penalty.}
A registered robustness check, descriptive by registration, refit the penalty \SOneSpecN{} ways for each rating. It varied the model type (for example a linear model, or an ordinal model that treats the five answers as ordered categories rather than evenly spaced numbers), its baselines (which differences between participants and between images the model allows for), its background variables (such as age and gender) and whether ratings were first averaged per person. In all \SOneSpecNTotal{} versions the AI label still lowered the rating. Among the versions that give the penalty in rating points (all but the ordinal ones), the penalty ran from \SOneSpecCreatMin{} to \SOneSpecCreatMax{} points on creativity and from \SOneSpecEffMin{} to \SOneSpecEffMax{} on effort. Descriptive by registration means that neither the middle of these estimates nor the share of versions with a significant penalty counts as evidence for a hypothesis.

The other registered checks agreed. The ordinal model on its own gave a penalty on both ratings ($p$~\SOneOrdCreatP{} on creativity and $p$~\SOneOrdEffP{} on effort, where a $p$-value is the chance of a penalty at least this large if the label had no effect). Its estimates are on another scale, so their size cannot be compared with the rating-point penalties. Adding age, gender and formal artistic training as background variables gave penalties of \SOneCovCreat{} and \SOneCovEff{} points. Refitting the model once without each participant, and once without each image, moved the penalty by at most \SOneInfluenceCreatImage{} points on creativity and \SOneInfluenceEffPerson{} on effort, and never changed its sign. One effort item, on the first image with a human figure, labeled its top two answers differently from the others; dropping it moved the effort penalty by \SOneItemHOneEffShift{} points and the difference between the effort and creativity penalties (Study~1's H2) by \SOneItemHTwoShift. That last check was settled after the estimates had been read. The two items also differed in wording. Only the creativity item carried a definition (`the presence of at least a minimal original, intellectual contribution to the creation of the work'), and it pointed raters toward an intellectual contribution, which both descriptions gave to Alex (`Alex had a new idea'); the effort item carried none. Part of the difference between the two penalties may therefore come from the wording of the items.

\paragraph*{The belief check.}
Every image was AI-made, so the Human label was a cover story, and no question asked participants whether they believed it. Instead, the registered belief check counts what Human-label participants wrote in the short explanation required for every rating, using fixed word lists that a script applied without seeing the label or the rating. Of the \SOneBeliefNHuman{} Human-label participants, \SOneBeliefBroadN{} (\SOneBeliefBroadPct; 95\% confidence interval, CI, \SOneBeliefBroadLoPct{} to \SOneBeliefBroadHiPct) mentioned machines, prompting or generating at least once, and \SOneBeliefNarrowN{} (\SOneBeliefNarrowPct; 95\% CI \SOneBeliefNarrowLoPct{} to \SOneBeliefNarrowHiPct) wrote at least once that a work is or looks machine-made. A 95\% CI is the range of values the data do not rule out at the 95\% level; these come from drawing participants again at random 2,000 times. Both are counts of statements, not of disbelief. A participant who doubted the label may have written nothing about it, and one who wrote that a work looks AI-generated may still have accepted the description. The rule that flags a statement that a work is or looks machine-made was settled in September and October 2026, after results by label had been read.

Disbelief in the cover story would make Human-label ratings more like AI-label ones, shrinking the penalty, not creating it. In a registered refit without those \SOneBeliefNarrowN{} participants, the penalty grew to \SOneBeliefRefitCreatDiff{} points on creativity and \SOneBeliefRefitEffDiff{} on effort. The refit is a sensitivity check and does not replace the full-sample estimate.

Dropping low raters from the Human-label group enlarges a penalty by arithmetic alone. To gauge that, an analysis that the registration names as exploratory description dropped 2,000 random groups of \SOneBeliefNarrowN{} Human-label participants whose mean effort rating matched the excluded group's within a fixed tolerance. It compared Cohen's $d$, the penalty divided by the pooled standard deviation (a measure of spread) of participants' mean ratings, before and after each removal. On creativity, $d$ rose from \SOneDCreat{} in the full sample to \SOneBeliefRefitCreatD{} after the exclusion, inside the middle 95\% of the random removals (\SOnePlaceboCreatNullLo{} to \SOnePlaceboCreatNullHi), so the change is what removing any such group produces. On effort, $d$ rose from \SOneDEff{} to \SOneBeliefRefitEffD, above every random removal (middle 95\%: \SOnePlaceboEffNullLo{} to \SOnePlaceboEffNullHi). The registered reading of the effort result is only that the change is not compatible with that reference. The matching held the mean effort rating alone, so the comparison cannot say which other difference between the groups moved the estimate.
